\paragraph{Access as a condition, not a guarantee.} By \cref{eq:scaling}, an edit to a page cannot act through a
first-party route that the engine cannot use. Study~1 shows how often sites present potential obstacles on such routes:
\BarAny\% of the organisational sites showed at least one (\BarNoEdge\% without research-client edge refusals), and for
most of them none of the flagged conditions was expressed through an agent-specific rule, so an audit of AI-specific
rules alone would miss them. These are measured potential obstacles, not exclusion: robots.txt records policy, our edge
requests came from a declared research client rather than a verified crawler, and the JavaScript measure is a
word-count proxy. The components can overstate and understate practical access in different directions, so the
composite is not a bound on exclusion. An obstacle also affects only the route through the organisation's own site;
third-party pages about it remain readable, which is one reason earned coverage mattered in Study~2.

\paragraph{What the audits and the search task add.} The live audit qualifies access in two ways. A fifth
(\AudPBlockedPct\%) of Perplexity's displayed citations went to pages whose robots.txt disallows its crawler, so a
disallow rule did not keep those pages out of its answers; a page can also reach an engine through a partner index, an
earlier crawl or a user-requested fetch. Access is a condition for an engine to use a site's current content through
its own crawler, not control over being cited. And ChatGPT's free tier, the most used assistant's default, showed no
visible search or citations in our \AudCNQ{} sessions while still recommending products. The interface does not reveal
how those answers were produced; if no retrieval occurred, they drew on what earlier crawls and training captured, and
sites restrict training crawlers far more often than search access (Study~1). Blocking a training crawler closes that
operator's documented collection route for an organisation's own pages; it does not guarantee that the pages are absent
from training data, and third-party pages about the organisation may still be used. The search task shows a related limit inside an engine that does search: the models opened
the target's page in only \AgBaseOpen\% of baseline answers, and the citation, statistics and instruction packages,
which sit below the page's opening paragraph, had no detectable effect. That task also changes the candidate pool,
snippets and tools, so it does not isolate retrieval as the cause, but it shows that effects measured with pages
supplied need not survive when models choose what to read. How often production assistants retrieve, and which pages
they open, limits how far any content-editing result carries over, ours included.

\paragraph{Endorsement and evidence.} The strongest lever was a recommendation from a third party, in line with the
weight that AI search gives earned media \citep{chen2025dominate}. It raised mentions significantly in every main model
and again had the largest pooled effect in the replication, and in an exploratory follow-up being singled out added
about as much again as being mentioned (\cref{sec:explore}). The package combines a third-party source with new
information and an explicit recommendation. A follow-up that held the words fixed on the brand's own page found that
recommending wording mattered more than the stated source (\AttrRecU{} against \AttrExtU{} points;
\cref{sec:attribution}); it did not test whether the same words weigh more on a separate site. Among edits
to the brand's own page, the citation package was the strongest, and statistics, citations and a quotation all beat the
unedited page, as \citet{aggarwal2024geo} reported. Unlike them, we find that keyword stuffing helped about as much as
statistics and persuasive copy also helped, so models did not reward evidence alone. What set the evidence packages
apart was what the answer said: by a keyword detector, answers that named the brand repeated the added sources and
figures far more often than persuasive phrases. Restating facts as questions, a summary or a table did not change
mentions significantly, and freshness dates and stated limitations did no better than the placebo. The summary did
help a page get found: in the search task it put the page among the results \AgESurfAnswerFirstU{} points more often and
raised mentions by \AgEMenAnswerFirstU{} points. Endorsement remained effective in search and was the only intervention
that raised the chance of being named first there.

\paragraph{Long pages and search.} On pages about six times longer, one sentence of evidence or praise became a small
share of what the model read, and its effect shrank or vanished, while a sentence repeating the question's words
gained. One possible explanation, which we did not test, is that models use overlap with the question to decide which
part of a long page is relevant, much as position and length shape what they use from long contexts
\citep{liu2024lost}. The search task points the same way: in our retriever, text that matched the question made a page
easier to find, and the gain carried through to mentions. The practical reading is not to stuff keywords, which search
engines penalise \citep{google2026spam}, but to state plainly, near the top of the page, which questions and needs the
offering answers, in the words buyers use, and to keep evidence close to that statement. Our retriever, unlike a search
engine, does not demote keyword stuffing, so the summary is the form of this advice more likely to transfer.
Endorsement weakened on long pages too, possibly because the recommendation became one sentence in a longer guide;
prominence within third-party coverage may matter as well as its existence.

\paragraph{Ranking and competition.} Answer engines retrieve through search indexes, their own or a partner's, so by
\cref{eq:scaling} a page must rank for the sub-queries an engine issues before its content can matter. Study~2 held this
stage fixed by supplying the pages. Answer engine optimisation therefore adds to search engine optimisation, chiefly by
requiring access for a new set of crawlers, rather than replacing it. Gains were relative: treatments raised the
target's share of brand mentions while the number of brands named barely changed, so the target displaced
competitors. If competitors made the same edits the relative gains could shrink, although the benchmark did not test
this. Statistics, citations and keywords are easy to copy; a third party's recommendation and a high search ranking are
not.

\paragraph{Manipulation.} A one-sentence instruction embedded in the page text was followed by most models and almost
never disclosed, while Anthropic's models moved against it or called it out, so resistance is achievable, although our
data cannot say how it was achieved. Hidden text has steered a commercial AI search product
\citep{guardian2024chatgpt}, a form of indirect prompt injection \citep{greshake2023indirect}. In the search task the
instruction had no detectable effect, but the models seldom opened the page that carried it; agents that read pages in
full may remain exposed. Any system that passes retrieved text to a model without treating it as data, including
on-device assistants and browsing agents, risks inheriting this weakness. Operators should treat text in retrieved
pages as data, remove hidden text before generation and test whether their models follow instructions found in
sources; we did not evaluate any of these mitigations. Measures based on the first brand named must also set aside
answers that name a brand only to warn about it. Organisations should not use the technique: it deceives users,
violates search engines' spam policies \citep{google2026spam} and can be detected in the page source.

\paragraph{Implications.} For organisations, the evidence, summarised technique by technique in \cref{app:evidence}, is
consistent with the following order of work, although no step was tested for its effect on visibility in a deployed
engine:
\begin{enumerate}
  \item Check access for each engine: evaluate robots.txt for each answer agent's token, including what the
  \texttt{*} group excludes; make sure robots.txt never returns a server error; ask the content-delivery network
  whether verified AI crawlers are admitted; and serve the main content in the HTML.
  \item Decide training and answering separately. Blocking a training crawler does not require blocking the same
  operator's search crawler, and \PctOpenAISplit\% of the sites that block \agent{GPTBot} already keep
  \agent{OAI-SearchBot} open.
  \item Earn coverage from independent sources and rank in the search indexes that engines use.
  \item Put checkable evidence, such as figures, sources and named evaluations, on the pages engines retrieve; it is
  what answers repeated when the page was in context.
  \item Say near the top of each page which questions and needs it answers, in the words buyers use. On longer
  synthetic pages this mattered as much as endorsement, and in the search task a summary at the top made the page
  easier to find. Keyword stuffing achieves the same crudely, is easy to copy and violates search engines' spam
  policies.
  \item Do not rely on restating content as questions and answers or on freshness stamps, which did nothing here, or on
  superlatives, which stopped working on longer pages; never hide instructions in a page.
\end{enumerate}
For operators and standards bodies, the results argue for separate tokens for training, search indexing and
user-requested fetches at every operator; published verification, such as IP ranges or signed requests, so that
networks can admit answer agents without admitting impostors; and documentation of how each crawler treats groups that
do not name it.
